% Generated draft: requires mathematical and source review.
\documentclass[11pt]{article}
\usepackage[margin=25mm]{geometry}
\usepackage{fontspec}
\setmainfont{Latin Modern Roman}
\usepackage{amsmath,amssymb,mathtools,unicode-math}
\setmathfont{Latin Modern Math}
\usepackage{enumitem,needspace,xurl,hyperref}
\setlength{\parindent}{0pt}
\setlength{\parskip}{5pt}
\setlength{\emergencystretch}{3em}
\newcommand{\N}{\mathbb{N}}
\newcommand{\Z}{\mathbb{Z}}
\newcommand{\Q}{\mathbb{Q}}
\newcommand{\R}{\mathbb{R}}
\newcommand{\C}{\mathbb{C}}
\newcommand{\F}{\mathbb{F}}
\newcommand{\A}{\mathbb{A}}
\newcommand{\PP}{\mathbb{P}}
\newcommand{\E}{\mathbb{E}}
\newcommand{\eps}{\varepsilon}
\newcommand{\dd}{\,\mathrm{d}}
\begin{document}
\section*{UnsolvedMath ID: 8400038}
\textbf{Source problem code:} catalogue entry
\subsection{Definitions and mathematical statement}
The source asks whether the ``same binary quadratic Diophantine solvability problem'' is NP-hard under randomized reductions. Its antecedent problem is not supplied. A binary quadratic input might specify an integer polynomial in two variables, but the exact equation, coefficient restrictions, variable domain, and size constraints are essential; integer, nonnegative-integer, and bounded-solution versions cannot be silently identified.

Once the intended decision language $L$ is fixed, NP-hardness under polynomial-time randomized many-one reductions in one common bounded-error sense means: for every $A\in\mathrm{NP}$ there is a randomized polynomial-time map $R$ such that for every input $x$,
\[
\Pr_r\bigl[(x\in A)\Longleftrightarrow(R(x;r)\in L)\bigr]\geq\frac23.
\]
The source does not specify many-one versus oracle reductions or bounded-error versus zero-error conventions, so this is explanatory rather than a verified exact reduction model. A deterministic polynomial-time many-one reduction is a special case of the displayed randomized model.

The supplied summary attributes a deterministic NP-completeness result to Manders--Adleman (1978) and says it settles this item. That implication requires verifying that their problem is exactly the missing antecedent language; the record does not provide enough information to do so. No particular quadratic equation is invented to fill the gap.
\subsection{Short English statement}
Is the source's binary quadratic Diophantine decision problem NP-hard using randomized reductions? The exact solvability problem and reduction convention are missing, so the supplied deterministic-hardness attribution cannot be matched without the preceding source item.
\subsection{Sources}
\begin{itemize}
\item Leonard Adleman and Kevin McCurley, \emph{Open Problems in Number Theory Complexity II} (1994), O33b, source page 22 / LNCS page 312, \url{https://mccurley.org/papers/open.ps.gz}; accessed 2026-07-29. The provenance records PDF text recovered against the coauthor-hosted PostScript.
\item Manders--Adleman (1978), as the deterministic NP-completeness attribution in the supplied research summary; no theorem statement or complete publication citation is given.
\item Local registry \texttt{data/unsolvedmath/problems.json}, permanent UnsolvedMath ID $8400038$, source code catalogue entry, \emph{O33b --- Randomized NP-hardness of binary quadratic Diophantine solvability}. Dataset attribution: Ulam AI and the UnsolvedMath Contributors.
\end{itemize}
\end{document}
